\documentclass[11pt]{article}
\usepackage[utf8]{inputenc}
\usepackage[T1]{fontenc}
\usepackage{amsmath,amssymb}
\usepackage[margin=1.15in]{geometry}
\usepackage{booktabs}
\usepackage{authblk}
\usepackage[colorlinks=true,linkcolor=blue,citecolor=blue,urlcolor=blue]{hyperref}

\title{The Tetrahedral Stereocentre from Packing:\\
\large Chirality's Geometric Origin without Hybridisation}

\author[1]{Claes Johnson}
\affil[1]{KTH Royal Institute of Technology, SE-100 44 Stockholm, Sweden\\
\small\texttt{claesjohnson@gmail.com}}
\date{\today}

\begin{document}
\maketitle

\begin{abstract}
Molecular chirality rests on one geometric fact: a carbon centre is tetrahedral, so four distinct
substituents admit two arrangements that are mirror images and not superimposable. That fact was
inferred from optical activity in 1874, half a century before quantum mechanics, and the orbital
account recovers it by \emph{hybridising} one $s$ and three $p$ eigenmodes into four equivalent
lobes --- a construction laid over an independent-particle approximation, undertaken because the
bare eigenmodes do not give molecular shape. We note that RealQM, a model of non-overlapping unit
electron charge densities in three space dimensions, obtains the tetrahedron directly: four unit
charges on a shell sit as far apart as the shell allows, and for $N=4$ that is tetrahedral, with a
Coulomb energy of $3.674/R$ against $6.000/R$ for the same four charges smeared over the shell.
The stereocentre is then a packing result rather than an orbital repair, and the existence of
chirality follows from how charge packs around a kernel. We are equally explicit about what this
does not deliver. Enantiomers come out exactly degenerate, which is correct but automatic in any
parity-symmetric theory and therefore buys nothing. Enantioselectivity does not come out at all:
a $95{:}5$ enantiomeric ratio corresponds to $\Delta\Delta G^{\ddagger} \approx 1.7$~kcal/mol,
three orders of magnitude below this model's accuracy floor. Chiral recognition energies are
typically under $1$~kcal/mol and dispersion-dominated, which is a permanent gap. And the origin of
biological homochirality requires parity violation, which the framework does not contain. The
contribution is therefore to chirality's \emph{geometric origin}, not to chiral chemistry.
\end{abstract}

\section{One geometric fact}

Chirality in carbon chemistry reduces to a single statement about shape. A centre bonded to four
groups is tetrahedral; if the four groups differ, there are exactly two ways to arrange them, and
the two are mirror images that cannot be brought into coincidence by any rotation. Everything
else --- optical rotation, enantiomeric excess, the handedness of sugars and amino acids --- is
built on that.

The historical order is worth keeping in view, because it is the reverse of the usual telling.
Van't Hoff and Le Bel proposed the tetrahedral carbon in 1874 \emph{in order to explain optical
activity}. The tetrahedron was an empirical inference from chirality, not a prediction about it,
and it preceded Schr\"odinger's equation by fifty-two years. Quantum mechanics inherited the
tetrahedron as something to be recovered.

\section{How the orbital account recovers it}

The one-electron hydrogenic states are $s$, $p$, $d$, $f$, classified by angular momentum. For
carbon the occupied valence modes are one $2s$ and three $2p$, and the $2p$ lobes point along three
mutually orthogonal axes. That is not a tetrahedron; orthogonal axes give $90^\circ$ between
bonds, where methane has $109.47^\circ$. The repair is \emph{hybridisation}: one forms four linear
combinations
\[
\tfrac12\left(s + p_x + p_y + p_z\right), \quad
\tfrac12\left(s + p_x - p_y - p_z\right), \quad \text{and two more},
\]
each pointing at a cube corner, and the four together are tetrahedral. VSEPR does the same work by
a different route, positing that electron pairs repel and arrange themselves at maximum separation.

Neither device is in the Schr\"odinger equation. The $N$-electron equation has no orbitals at all;
its solutions are functions on $3N$ variables carrying no $s,p,d,f$ labels. The orbital scheme is
an independent-particle approximation laid over it, and hybridisation is a further construction
laid over that, undertaken --- as the companion article on atoms puts it --- because ``the bare
eigenmodes do not give molecular shape.'' VSEPR and $sp^3$ hybridisation are geometric repairs
grafted onto the orbital model precisely to recover a shape the decomposition had destroyed.

This is not an objection to the orbital account, which is enormously useful and which we do not
dispute. It is an observation about where the tetrahedron comes from in that account: it is put
back in, by hand, after an angular decomposition removed it.

\section{How packing gives it directly}

RealQM carries no angular decomposition to undo. The state is a set of non-overlapping unit charge
densities, and the question of where four charges sit on a shell is answered by minimising Coulomb
repulsion subject to non-overlap. The answer is the obvious one: as far apart as the shell allows
--- antipodal for two, trigonal for three, \textbf{tetrahedral for four}.

The numbers make the point quantitative. For unit charges at a common radius $R$, in units of
$1/R$:

\begin{center}
\begin{tabular}{ccc}
\toprule
charges & smeared over the shell & packed \\
\midrule
2 & 1.000 & 0.500 \ \ (antipodal) \\
3 & 3.000 & 1.732 \ \ (trigonal) \\
4 & 6.000 & \textbf{3.674} \ \ (tetrahedral) \\
\bottomrule
\end{tabular}
\end{center}

Packing four unit charges tetrahedrally costs $3.674/R$ where smearing the same four charges over
the shell costs $6.000/R$ --- a factor $1.63$. The tetrahedron is not imposed to recover a known
shape; it is what four mutually repelling non-overlapping densities do. The same rule gives the
shell capacities of the periodic table, where the closed octet appears as $8 = 4{+}4$, two dual
tetrahedra, and where there is no $s$, $p$ or $d$ anywhere in the construction.

\paragraph{So the stereocentre is a packing result.} Chirality exists because four unit charges
pack tetrahedrally around a kernel, and a tetrahedron with four distinct corners is chiral. No
hybridisation is required because nothing was decomposed into orbitals in the first place. Of the
two classical repairs, VSEPR is much the closer in spirit --- it is a repulsion-and-separation
argument, and RealQM is in effect VSEPR promoted from a rule to a solved variational problem.

\section{What this does not deliver}

The above is a claim about the \emph{origin} of chirality and nothing more. Four limits should be
stated plainly, because each of them excludes something a reader might reasonably hope for.

\paragraph{Enantiomer degeneracy is automatic, and therefore worthless as evidence.} Two
enantiomers have exactly equal energy in any parity-symmetric theory. RealQM is Coulomb-only and
hence parity-symmetric, so the degeneracy comes out exactly --- and so it does in standard quantum
mechanics. Nothing is learned either way.

\paragraph{Enantioselectivity is out of reach by three orders of magnitude.} Selectivity is a
transition-state energy difference. A $95{:}5$ enantiomeric ratio at room temperature corresponds
to $\Delta\Delta G^{\ddagger} \approx 1.7$~kcal/mol. This model's energies are accurate to a few
per cent of a total energy, which for an organic molecule is hundreds of kcal/mol. The gap is not
a matter of refinement.

\paragraph{Chiral recognition energies lie in the dispersion gap.} Discrimination between
enantiomers by a chiral host is typically under $1$~kcal/mol and is substantially dispersive.
RealQM's densities are static, so it has no dispersion; the same limitation that keeps it out of
van der Waals complexes keeps it out of recognition thermodynamics. What it \emph{can} address is
the shape question --- whether an enantiomer geometrically fits a pocket --- since hydrogen-bond
and steric geometries come out at measured distances. Whether it fits, not how much it prefers to.

\paragraph{Homochirality requires physics the framework lacks.} That life uses L-amino acids and
D-sugars is a broken symmetry, and breaking it requires parity violation --- the weak interaction.
RealQM contains Coulomb interaction and nothing else. The origin of biological handedness is
outside it, and no amount of computation changes that.

\section{What is worth saying, then}

Exactly one thing, and it is upstream of chiral chemistry rather than part of it. The founding
geometric fact of stereochemistry --- that a four-coordinate centre is tetrahedral --- is in the
orbital account a repair applied after an angular decomposition, and in RealQM a direct consequence
of four non-overlapping charges minimising their mutual repulsion. Chirality is then a packing
phenomenon: it exists because charge packs tetrahedrally, and a tetrahedron with four distinct
corners has no mirror symmetry.

Van't Hoff and Le Bel inferred the tetrahedron from optical activity in 1874. Pauling supplied
hybridisation in 1931 to recover it from the new mechanics. The claim here is that a model built
from charge densities in three dimensions does not need the second step, because it never takes the
decomposition that made it necessary.

\end{document}
